\section{Thermal stability of a retuned zero}
\label{sec:thermal}

\subsection{An operator zero and a ground-state zero}

A vanishing ground-state source does not establish that a probe vanishes as
an operator. Temperature distinguishes these possibilities by populating
additional initial states. We keep $\tau=T/J_{\rm eff}$ fixed as
$\beta\to0$, where $J_{\rm eff}=\beta c_1(\alpha_*)^2$ and $k_B=1$.
Thus the thermal ensemble remains inside the local ground multiplets.
Put $h=\sum_i\bm T_i\cdot\bm T_{i+1}$ and let $\epsilon_i$ be its
dimensionless eigenvalues. For the Hermitian leading $Q=\pi$ probe,
write $D=D_\pi$, $K=K_q$, and define
\begin{align}
 R_\kappa&=K+\kappa A'D,\notag\\
 C_\kappa(\tau)&=\sum_{\epsilon_f>\epsilon_i}
 p_i\left|\langle f|R_\kappa|i\rangle\right|^2,
 \qquad p_i=\frac{e^{-\epsilon_i/\tau}}{Z}.
 \label{eq:thermal-weight}
\end{align}
This coefficient describes the regular positive-frequency response on the
$J_{\rm eff}$ scale. At fixed finite size, choose $\eta>0$ smaller than
the smallest nonzero level spacing of $h$, and let $W_{+,\eta}$ include
only low-band transitions with $\omega>\eta J_{\rm eff}$. Then
$W_{+,\eta}=\beta^2C_\kappa+O(\beta^3)$. This qualification matters:
higher-order terms can split degenerate levels of $h$ and move weight
from zero frequency to $\omega=O(\beta^2)$. That quasielastic weight is
not automatically included in $C_\kappa$, even though it contributes to
a microscopic integral over every $\omega>0$.
All equal-energy blocks, including off-diagonal elements within degenerate
subspaces, are excluded. A real exchange retuning is optimized separately
at each temperature, so a mere displacement of the zero cannot account for
a nonzero minimum.

At any $\tau>0$ every state of a finite effective chain has positive
Boltzmann weight. Nonnegativity of the summands and Hermiticity imply
\begin{equation}
 C_\kappa(\tau)=0\quad\Longleftrightarrow\quad[h,R_\kappa]=0.
 \label{eq:thermal-criterion}
\end{equation}
The right-hand side allows an operator acting within equal-energy blocks;
it need not be the zero operator. This standard positivity observation is
useful here because it replaces a temperature-by-temperature cancellation
test with a local operator calculation. For a doublet, the stronger identity
$K=\gamma D$ already proved above cancels $R_\kappa$ throughout the
retained space. Consequently $W_+=O(\beta^4)$ or smaller at every fixed
$\tau$ under the same periodic and reflection assumptions.

\subsection{A representation-dependent obstruction}

For the quartet root at $q=\pi$, define the unnormalized alternating spin
$\widetilde D=\sum_i(-1)^iT_i^z$ and
$Y=\sum_i(-1)^iX_{i,i+1}$, using Eq.~\eqref{eq:bond-identity}. Then
\begin{equation}
 R_\kappa=\frac{bY+(-2a+\kappa A')\widetilde D}{\sqrt{5N}}.
 \label{eq:thermal-source}
\end{equation}
Both $\widetilde D$ and $Y$ are spin vectors with momentum $\pi$,
odd under time reversal and even under the site-permutation reflection
$i\mapsto-i$. For $Y$, reflection reverses the bond orientation and
shifts its alternating sign, leaving their product unchanged.
The two contributions therefore share these symmetry labels. The issue is
not whether a transition is symmetry-allowed, but whether one dipole
coefficient can cancel all its dynamical amplitudes across different
initial and final states.
This is consistent with the independent reduced matrix elements in the
multiplicity case of Ref.~\onlinecite{tabrizi2021selection}, Eq.~(16).
Their independence as operator structures alone is not a proof of
noncancellation: the inelastic matrix elements could still be
proportional, with their difference acting only inside equal-energy
blocks. The commutator minimum below excludes precisely this possibility
for the specified effective chain.
Different bond-spin channels obstruct cancellation of its commutator with
$h$. The spectral-moment approach is standard in spin-chain
dynamics~\cite{mueller1982sumrules}; here it gives an exact retuning
minimum for this dipole--composite probe. More generally, take a spin-$s$
chain with the same uniform nearest-neighbor
$h$ and probe $R=(bY+u\widetilde D)/\sqrt{\ell N}$, where
$\ell>0$ fixes the probe normalization ($\ell=5$ for the pentamer).
Let $d=2s+1$ and $c=s(s+1)$. For every finite even $N\ge4$, the
uniformly traced, two-sided second frequency moment is exactly
\begin{align}
 M_2(u)&=\frac{1}{d^N}\operatorname{Tr}[h,R]^\dagger[h,R]\notag\\
 &=M_{2,\min}+\frac{(d^2-1)^2}{18\ell}
       \left[u+\frac{b(d^2-4)}{12}\right]^2,\\
 M_{2,\min}&=\frac{b^2c^2(4c-3)(86c-27)}{810\ell}.
 \label{eq:general-moment}
\end{align}
Appendix~\ref{app:general-moment} derives this identity from local spin
traces and the orthogonality of distinct operator supports. It is not a
fit to a finite set of spins or chain lengths. Independent full-chain
matrices for $s=1/2,1,3/2$ at $N=4,6$ reproduce the identity.

For $s=1/2$ the quadrupole vanishes and the minimum is zero. For
$s\ge1$ and $b\ne0$, it is strictly positive. The conclusion is
conditional on this probe form and on retuning only its dipole
coefficient: a higher local spin alone does not guarantee a nonzero $b$
in an arbitrary microscopic cluster. For the pentamer quartet the exact
minimum is $(197/16)b^2$. Combining this positive minimum with
Eq.~\eqref{eq:thermal-criterion} proves $C_{\min}(\tau)>0$ for every
finite even $N\ge4$ and every $\tau>0$ in the leading quartet model.
A ground-state cancellation therefore cannot be continued into complete
thermal cancellation by the same one-parameter retuning.
The moment optimum $u=-b$ differs from the four-cell ground-state optimum
$u=3b$; minimizing one observable need not minimize another.

This trace includes positive and negative frequencies; its positive-only
uniform-trace value is half as large. It refers to the retained effective
space, not the full microscopic infinite-temperature spectrum. In the
fixed-size weak-coupling limit it describes the ordered regime
$J_{\rm eff}\ll T\ll\Dloc$. At the four-cell zero-temperature retuning,
the prediction is $M_2=(837/16)b^2=0.79932812$. Directly summing the
twenty-spin low-band transitions, with uniform initial weights and
division by $\beta^2J_{\rm eff}^2$, gives $0.79967945$ at
$g=0.000625$, a $0.044\%$ deviation; the deviation increases to
$0.088\%$ and $0.177\%$ at the next two couplings. No frequency window
is used for this check: the frequency-squared factor suppresses the
quasielastic contribution in the leading limit. This does not remove the
separate finite-momentum limitation or by itself yield an integrated-weight
lower bound.

Locality supplies a stronger conclusion specifically in the uniform-trace
limit. Let $C^{\rm tr}_{>\nu}$ denote the positive-frequency weight of
$R$ with dimensionless frequency above $\nu$, using equal initial-state
weights $1/d^N$. Appendix~\ref{app:weight-bound} proves, for the pentamer
quartet and every real $u$,
\begin{equation}
 C^{\rm tr}_{>\sqrt{45}/4}\ge\frac{197}{1036800}b^2>0
 \qquad(N\ge4\text{ even}).
 \label{eq:quartet-window-bound}
\end{equation}
This uses a fourth-moment upper bound as well as
Eq.~\eqref{eq:general-moment}. It excludes a sequence of larger even
rings whose optimized uniform-trace response vanishes or moves entirely
to arbitrarily small frequencies. The coefficient is conservative, not
a predicted intensity. The statement is within the leading effective
model with the uniform-trace limit taken before increasing $N$; it does
not establish a lower bound at arbitrary fixed $T/J_{\rm eff}$ or
uniform validity of the microscopic weak-coupling expansion with size.

\subsection{Independent local controls do not evade the obstruction}

The preceding result might appear to reflect an insufficient number of
control parameters. At the quartet root $q=\pi$, however, any bounded
SU(2)-invariant perturbation of order $\beta$ confined to an individual
cell changes its projected vector probe only by a multiple of $T_i^z$.
The isolated cell contains a single spin irrep; its projected scalar
perturbation shifts all components equally. The leading uniform $h$ and
the coefficient $b$ are therefore unchanged. Even independent local
controls in every cell can only enlarge the leading probe to
\begin{equation}
 R_{\{u_i\}}=\frac{bY+\sum_i(-1)^iu_iT_i^z}{\sqrt{\ell N}},
 \qquad u_i\in\mathbb R.
 \label{eq:local-control-probe}
\end{equation}
This statement concerns infinitesimal local retunings, not arbitrary
changes to the leading intercell exchange or to the local spin multiplet.

Translation symmetry separates the $Q=\pi$ component, whose dipole
coefficient is $\bar u=N^{-1}\sum_i u_i$, from all other operator momenta.
Their thermal spectral weights add without interference. For either the
positive-frequency weight or any fixed positive-frequency window,
\begin{equation}
 F(R_{\{u_i\}})=F(R_{\bar u})+
 F(R_{\{u_i\}}-R_{\bar u})\ge F(R_{\bar u}).
 \label{eq:local-control-projection}
\end{equation}
Here $F$ uses a Gibbs state of the uniform leading $h$, or its uniform
trace. Thus the minimum over all real $u_i$ equals the uniform minimum.
Both the finite-temperature obstruction and the uniform-trace window
bound survive this larger control space. Appendix~\ref{app:local-controls}
gives the proof and microscopic tangent checks. The projection is an
operator identity used in the proof, not an averaging of plotted spectra.

\subsection{Exact activation from a dark four-cell ground state}

The four-cell ring supplies a solvable example of why the distinction
matters. Write $\bm A=\bm T_0+\bm T_2$ and
$\bm B=\bm T_1+\bm T_3$, so that $h=\bm A\cdot\bm B$.
Its ground state has $(j_A,j_B,J)=(3,3,0)$ and energy $-12$;
the first excited multiplet is $(3,3,1)$ at $-11$.
The zero-temperature optimum is $\lambda_0=-2a-3b$, with
$\lambda=-\kappa A'$. The resulting source
$b(Y+3\widetilde D)/\sqrt{20}$ annihilates the ground state exactly,
although it does not commute with $h$.

The first thermally occupied triplet opens two groups of transitions.
Their final sectors are $(3,3,2)$ at $\omega/J_{\rm eff}=2$ and
$(3,1,2),(1,3,2)$ at $\omega/J_{\rm eff}=7$.
Summed over the three initial magnetic components, the squared matrix
elements divided by $b^2$ are respectively $27/50$, $6/25$, and $6/25$.
Consequently,
\begin{equation}
 C_{\min}(\tau)=\frac{51}{50}b^2e^{-1/\tau}
                  +o(e^{-1/\tau}).
 \label{eq:thermal-activation}
\end{equation}
The two groups have leading intensity ratio $9:8$. Reoptimizing the
exchange at each temperature changes this result only beyond the leading
activation order. Exact rational matrix elements reproduce the independently
assembled full-space thermal calculation to relative error
$3\times10^{-15}$ in its leading coefficient. The activation gap, line
positions, and intensity ratio are four-cell results, not universal
long-chain spectral predictions.

\begin{figure*}[t]
\centering
\includegraphics[width=\textwidth]{thermal_evidence.pdf}
\caption{Thermal cancellation and its momentum-acceptance limitation.
(a) Minimum leading positive-frequency coefficient in four- and six-cell
effective periodic chains. The doublet zero is the operator identity,
not a numerical intensity cutoff. (b) Full microscopic two-cell weights
at $T/J_{\rm eff}=1.6$, normalized separately to their values at
$g=\beta/\Dloc=0.005$. Quartet and doublet slopes are $2.02549$ and
$4.00189$; this normalization does not equate their absolute intensities.
(c) Weights averaged over the uniform interval $q_0\pm\delta q$, relative
to their ideal central-momentum values, at $g=0.005$ and
$T/J_{\rm eff}=1.6$. The quartet and doublet central momenta are
$\pi$ and $3\pi/5$, respectively. Exchange parameters remain on their
central-momentum retuning trajectories. Both initial and final microscopic
states lie below $E_0+\Dloc/2$, while the partition function includes
the full spectrum. Lines connect computed points; no spectral or image
smoothing is applied. These are finite-chain results, not an instrument
resolution estimate.}
\label{fig:thermal-evidence}
\end{figure*}

\subsection{Microscopic check and finite momentum acceptance}

Full thermal traces on four- and six-cell effective chains give the
contrast in Fig.~\ref{fig:thermal-evidence}(a). At $\tau=1.6$, the quartet
minima are $0.00991078$ and $0.00968212$, respectively. For six cells the
latter is about 22 times its zero-temperature value, even after
temperature-dependent retuning. This finite-size comparison supports the
operator result without replacing it by a size extrapolation.

To test the perturbative prediction in the original model, we diagonalized
all 1024 states of two coupled pentamers. Two distinct microscopic closing
links require two corresponding effective exchange bonds.
The quartet root is Eq.~\eqref{eq:rank-exact-root}; the doublet uses
$q=3\pi/5$, $\alpha_*=1.09755143$, so both have $Q=\pi$.
At $g=0.005,0.01,0.02$ and $\tau=0.4,1.6$, the retuned quartet
weights have logarithmic slopes $2.07089,2.02549$, compared with
$4.00185,4.00189$ for the doublet. At the weakest coupling the quartet
coefficients differ from their predictions by $1.75\%$ and $1.18\%$.
At $\tau=0.4,g=0.02$ the error reaches $12.25\%$; quantitative
agreement is not claimed throughout the tested interval. These two-cell
calculations test the finite-temperature expansion and order contrast.

We tested the four-cell activation directly in the twenty-spin periodic
microscopic model at $g=0.000625,0.00125,0.0025$. All 256 retained states
were reconstructed by total-spin raising from the 44 low-band states in
the $S^z=0$ sector. The exchange was kept on the same leading
zero-temperature retuning trajectory, $\kappa_0=1.28029750$, rather than
reoptimized for each temperature. At the weakest coupling the first
thermal gap is $1.000318J_{\rm eff}$; the two line-group centroids are
$2.000592J_{\rm eff}$ and $7.000309J_{\rm eff}$. Their unweighted
transition strengths, summed over the initial triplet and divided by
$\beta^2b^2$, are $0.536037$ and $0.482098$, within $0.74\%$ of
$27/50$ and $12/25$. Both approach their predictions as $g$ decreases
[Fig.~\ref{fig:four-cell-activation}(b)].

The full thermal response also exposes two limits of the leading result.
At $\tau=0.1$ the weakest-coupling weight exceeds the prediction by
$25.9\%$: the ground-state residual accounts for $21.0\%$ of the
microscopic total and scales approximately as $\beta^{4.0015}$ over the
three tested couplings. It is retained in
Fig.~\ref{fig:four-cell-activation}(a), not subtracted to reveal an
activation law. At $\tau=1.6$, quasielastic lines originating from split
effective degeneracies contribute $1.68\%$ of the total positive-frequency
weight. A resolved window $\eta=0.05$ separates these lines from the
$O(J_{\rm eff})$ response. Its weakest-coupling coefficients agree with
the fixed-retuning effective predictions to $0.161\%$, $0.229\%$, and
$0.042\%$ at $\tau=0.2,0.4,1.6$; changing $\eta$ from $0.02$ to
$0.1$ leaves the computed window weights unchanged. The window was
introduced as a diagnostic after the threshold sensitivity was detected;
the original unwindowed results are retained. These results validate a
finite-frequency activation mechanism with a finite-coupling crossover,
not an activation law uniformly valid down to $T=0$.

\begin{figure*}[t]
\centering
\includegraphics[width=\textwidth]{four_cell_activation.pdf}
\caption{Direct four-cell microscopic check and low-temperature background.
(a) Resolved thermal weight for twenty physical spins with periodic
boundaries, $q=\pi$, fixed zero-temperature retuning, and
$\omega>0.05J_{\rm eff}$. Both initial and final states are in the
complete 256-state low band. The finite-coupling curves retain the
ground-state background. The effective curve is the full leading thermal
trace at the same retuning; the activation asymptote is
$(51/50)b^2\exp(-J_{\rm eff}/T)$.
(b) Strengths from the first thermally occupied triplet to the two predicted
final-state groups, divided by $\beta^2b^2(27/50)$ and
$\beta^2b^2(12/25)$, respectively. These strengths are summed over initial
magnetic components before thermal weighting. Lines connect calculated
points without smoothing. The omitted-band population has a rigorous
bound, and the frequency-window diagnostic is separate from the original
unwindowed calculation.}
\label{fig:four-cell-activation}
\end{figure*}

The ideal momentum condition is restrictive. We also averaged the
microscopic transition probabilities over a uniform window around each
central momentum, keeping the central retuning fixed. At
$g=0.005,\tau=1.6$, a half-width $0.001\pi$ raises the quartet weight
by $0.094\%$, but the doublet weight by a factor of approximately
$2.93\times10^6$ [Fig.~\ref{fig:thermal-evidence}(c)].
The local doublet form factor has $\partial_qA=1.69998$ at its root,
whereas the quartet root at $q=\pi$ is stationary in momentum.
A fixed window therefore adds an ordinary form-factor background that
can overwhelm the doublet's very small fourth-order central signal.
The observed doublet coupling slope falls to approximately $0.00475$
at that window width. The operator-level distinction at ideal momentum
remains valid, but by itself is not a resolution-independent experimental
diagnostic.

Optimizing the exchange for the entire window reduces, but does not
eliminate, this background. At the stationary quartet root, write
$x=q-\pi$, $A_{qq}=\partial_q^2 A(\alpha_*,\pi)$, and let $B$ be the
positive-frequency thermal weight of $D$. If $\kappa_0$ minimizes the
central response, the remaining source is thermally orthogonal to $D$.
The leading window-dependent dipole amplitude is
$\beta A'(\kappa-\kappa_0)+A_{qq}x^2/2$.
Retuning removes its mean, leaving the variance of $x^2$.
For a uniform half-width $\delta q=\rho\sqrt{\beta}$, at fixed finite
size and scaled temperature, this gives
\begin{align}
 \kappa_{\rm win}&=\kappa_0-\frac{A_{qq}\rho^2}{6A'}+o(1),\notag\\
 \frac{\overline W^{\rm opt}_{+,\eta}}{\beta^2}
 &\longrightarrow C_{\min}+\frac{A_{qq}^2B}{45}\rho^4.
 \label{eq:window-retuning}
\end{align}
Here the regular-frequency window is the one defined above. Keeping
$\kappa_0$ fixed instead gives $1/20$ in place of $1/45$ in the
background coefficient. Thus window retuning removes its mean shift,
not the variation of the form factor across the measured interval.
In the ten-spin ring at $\tau=1.6$, the optimized microscopic
coefficients for $\rho=0.25,0.5,1$ agree with this limit within
$0.30\%$ at $g=0.00125$; deviations increase with coupling.

The chosen doublet root has nonzero $A_q=\partial_q A$, so its leading
window background is instead $A_q^2B(\delta q)^2/3$. Keeping it at
the fourth-order scale requires $\delta q=O(\beta^2)$, whereas the
stationary quartet root permits $O(\sqrt{\beta})$ at its quadratic
scale. These different tolerances reflect the orders of the selected
momentum zeros, not a universal resolution advantage of higher spin.
Furthermore, optimizing beyond the first-order trajectory reduces the
two-cell doublet's central low-band response to numerical zero: only
one singlet--triplet reduced transition amplitude is involved there.
For $\delta q=\rho\beta^2$ its remaining window weight approaches
$\beta^4 A_q^2B\rho^2/3$, with a $0.055\%$ deviation at the weakest
tested coupling. Thus a measured fourth power alone need not identify
an intrinsic central-momentum residual. These shrinking-window tests
do not establish detectability at fixed instrumental resolution or
complete second-order cancellation on longer chains.
